\subsection{二次型。}

定义 2. --- 假设环 A 交换。称 E 到 A 中的映射 Q 为 E 上的二次型，如果

\begin{enumerate}
\def\labelenumi{\arabic{enumi})}
\item
  对 α ∈ A 与 x ∈ E 有 Q(αx) = α²Q(x)；
\item
  映射 \(\Phi : (x, y) \to \mathrm{Q}(x + y) - \mathrm{Q}(x) - \mathrm{Q}(y)\) （\(\mathbf{E} \times \mathbf{E}\) 到 A 中的）是一个双线性型。
\end{enumerate}

双线性型 \(\Phi\)（它必然是对称的）称为与 Q 相伴的双线性型。若 \(\Phi\) 非退化，则称 Q 非退化。

由于 Q(2x) = 4Q(x)，由 2) 立即得出

\[
\Phi (x, x) = 2 \mathrm{Q} (x).\tag{13}
\]

特别地，若 A 是特征为 2 的环，则型 \(\Phi\) 是交错型。

我们称 E 的两个元素（相应地两个子集）关于 Q 正交，如果它们关于相伴的双线性型 Φ 正交。

设 \((x_{i})_{i\in I}\) 是 E 的元素的一个族，\((a_{i})_{i\in I}\) 是 A 的元素的一个族，其中除有限多个外均为零。对指标 \(i\)（使 \(a_{i}\neq0\) 者）的个数作归纳，容易证明有

\[
\mathrm{Q} (\sum_ {i} a _ {i} x _ {i}) = \sum_ {i} a _ {i} ^ {2} \mathrm{Q} (x _ {i}) + \sum_ {\{i, j \}} a _ {i} a _ {j} \Phi (x _ {i}, x _ {j}),\tag{14}
\]

最后一个和式遍取 I 的二元子集。

对每个双线性型 \(f\) （在 \(\mathbf{E} \times \mathbf{E}\) 上），定义二次型 \(\mathbf{Q}\) 为 \(\mathrm{Q}(x) = f(x, x)\)；于是双线性型 \(\Phi\) （与 \(\mathbf{Q}\) 相伴的）由 \(\Phi(x, y) = f(x, y) + f(y, x)\)（对 \(x, y\)，在 \(\mathbf{E}\) 中）定义。此外，若假设标量 2 有逆 \(\frac{1}{2}\) （在 \(\mathbf{A}\) 中），则存在唯一的对称双线性型 \(f\) 使 \(\mathrm{Q}(x) = f(x, x)\)，即 \(f = \frac{1}{2}\Phi\)；\(f\) 关于组 \(S = (x_1, \ldots, x_n)\) 的判别式也称为 \(\mathbf{Q}\) 关于 \(S\) 的判别式。因此在这种情形，\(\mathbf{E}\) 上的二次型与对称双线性型之间存在一一对应（参见习题 6）。

在自由模的情形，我们还有如下结果：

命题 2. --- 假设 A 交换且 E 具有基 \((e_i)_{i \in I}\)。于是，对 E 上的每个二次型 Q，存在一个双线性型 \(f\) （E \(\times\) E 上的），使对一切 \(x \in E\) 有 Q(x) = f(x, x)。对每个标量族 \((b_{ij})_{(i,j) \in I \times I}\)，它满足 \(b_{ij} = b_{ji}\)（对 \((i, j) \in I \times I\)），则存在唯一的二次型 Q 使

\[
\mathrm{Q} (e _ {i}) = b _ {i i}, \quad \Phi (e _ {i}, e _ {j}) = b _ {i j} \text {   for   } i \neq j,\tag{15}
\]

其中 Φ 表示与 Q 相伴的双线性型；于是 Q 由下列公式给出

\[
\mathrm{Q} (\sum_ {i} a _ {i} e _ {i}) = \sum_ {\{i, j \}} b _ {i j} a _ {i} a _ {j},\tag{16}
\]

最后一个和式遍取 I 的具有一个或两个元素的子集 \(\{i, j\}\)。

事实上，由于公式 (16) 只是公式 (14) 的一个改写，满足 (15) 的二次型 Q 的唯一性得证。为证其存在性，首先注意存在 A 的元素的一个族（ \(b_{ij}'\) ），使 \(b_{ii}' = b_{ii}\) 且 \(b_{ij}' + b_{ji}' = b_{ij}\)（对 \(i \neq j\)）；例如，给 I 配备一个全序结构（集合论，第 III 章，§ 2，第 3 小节，定理 1），并令 \(b_{ij}' = b_{ij}\)（对 \(i<j\)），令 \(b_{ij}^{\prime}=0\)（对 \(i>j\)），便得到这样一个族。由于诸 \(e_{i}\) 组成 E 的基，存在一个双线性型 \(f\) （\(\mathbf{E}\times\mathbf{E}\) 上的）使 \(f(e_{i},e_{j})=b_{ij}^{\prime}\)；令 \(\mathrm{Q}^{\prime}(x)=f(x,x)\)，并以 \(\Phi^{\prime}\) 表示与二次型 \(\mathrm{Q}^{\prime}\) 相伴的双线性型，则得 \(\mathrm{Q}^{\prime}(e_{i})=b_{ii}\) 且 \(\Phi^{\prime}(e_{i},e_{j})=f(e_{i},e_{j})+f(e_{j},e_{i})=b_{ij}\)。这就证明了我们的第二个断言。至于第一个断言，它立即得出，因为据唯一性，若二次型 Q 满足 (15)，则有 \(\mathrm{Q}(x)=\mathrm{Q}^{\prime}(x)=f(x,x)\)。

配备了由二次型 Q 所定义的结构的模 E 称为二次模。二次模 (E, Q) 到二次模 (E', Q') 中的同态是指一个 E 到 E' 中的线性映射 u，使 Q = Q' ∘ u；若 Φ 与 Φ' 是与 Q 和 Q' 相伴的双线性型，则对 x ∈ E, y ∈ E 有 Φ(x, y) = Φ'(u(x), u(y))；换言之，Φ' 是 Φ 在 u 下的原像（§ 1，第 1 小节）。两个 A-模 E 与 E' 上的两个二次型 Q 与 Q' 称为等价的，如果相应的二次模同构。

设 \((\mathbf{E}_{i}, \mathbf{Q}_{i})_{i\in\mathbb{I}}\) 是二次模的一个族，E 是诸模 \(E_{i}\) 的直和。称二次模 \((\mathbf{E}_{i}, \mathbf{Q}_{i})\) 的外直和为这样一个二次模：它是在 E 上配备由 \(\mathrm{Q}(\sum_{i}x_{i}) = \sum_{i}\mathrm{Q}_{i}(x_{i})\)（对 \(x_{i} \in E_{i}\)）定义的二次型 Q 而得到的。也称二次型 Q 为诸二次型 \(Q_{i}\) 的外直和。

若诸型 \(Q_{i}\) 非退化，则 Q 也非退化。

设 Q 是 A-模 E 上的一个二次型；若 F 是 E 的一个子模，且 Q 在模 F 的每个类上取常值，则由 Q 取商得到的 E/F 到 A 中的映射 \(\overline{Q}\) 显然是一个二次型，且 E 到 E/F 上的典范映射是二次模结构的同态。要使 Q 在模 F 的每个类上取常值，必须且只需 \(Q(x + y) = Q(x)\)（对 \(x \in E\) 与 \(y \in F\)），即（以 \(\Phi\) 表示与 Q 相伴的双线性型）\(Q(y) + \Phi(x, y) = 0\)（对 \(y \in F\) 与 \(x \in E\)）。取 x = 0，可见 \(Q(y) = 0\)（对 \(y \in F\)），从而 \(\Phi(x,y)=0\)（对 \(x\in E\) 与 \(y\in F\)）。换言之，若把二次模 (E, Q) 的核定义为 E 中满足 \(Q(x)=0\) 且 \(\Phi(x,z)=0\)（对一切 \(z\in E\)）的元素 x 的集合 N，则要使 Q 在模 F 的每个类上取常值，必须且只需 F 含于 (E, Q) 的核 N。不难验证 N 是 E 的一个子模。因此，要使 Q 在模 F 的每个类上取常值，必须且只需 F 由 N 的元素生成。

立即可见二次模 E/N 的核是 \{0\}。

命题 3. --- 设 h 是 A 到交换环 A' 中的一个同态。对 A-模 E 上的每个二次型 Q，存在唯一的 A'-模 A'⊗A E（第 III 章，第 2 版，附录 II，第 10 小节）上的二次型 Q'，使得

\[
\mathrm{Q} ^ {\prime} (1 \otimes x) = h (\mathrm{Q} (x))\tag{17}
\]

对一切 \(x \in \mathbf{E}\)。此外，与 Q' 相伴的双线性型 \(\Phi'\) 是由与 Q 相伴的双线性型 \(\Phi\) 经标量环扩张得到的。

我们先证明，若存在满足 (17) 的二次型 Q'，则它是唯一的，且与 Q' 相伴的双线性型 \(\Phi'\) 是由与 Q 相伴的型 \(\Phi\) 经标量环扩张得到的。事实上，后一断言由以下事实得出

\[
\begin{array}{r l} \Phi^ {\prime} (1 \otimes x, 1 \otimes y) & = Q ^ {\prime} (1 \otimes x + 1 \otimes y) - Q ^ {\prime} (1 \otimes x) - Q ^ {\prime} (1 \otimes y) \\ & = h (\Phi (x, y)) \end{array}\tag{18}
\]

对 \(x \in \mathrm{E}, y \in \mathrm{E}\)。公式 (14) 于是表明有

\[
\mathrm{Q} ^ {\prime} \left(\sum_ {i} a _ {i} ^ {\prime} \otimes x _ {i}\right) = \sum_ {i} a _ {i} ^ {\prime 2} h (\mathrm{Q} (x _ {i})) + \sum_ {\{i, j \}} a _ {i} ^ {\prime} a _ {j} ^ {\prime} h (\Phi (x _ {i}, x _ {j}))\tag{19}
\]

对 \(a_{i}^{\prime}\in\mathbf{A}^{\prime}\) 与 \(x_{i}\in\mathbf{E}\)，这就证明了 \(\mathbf{Q}^{\prime}\) 的唯一性。

为证 Q' 的存在性，先假设模 E 具有基 \((e_{i})_{i\in I}\)。于是存在 A 的元素 \(b_{ij}\)，使 \(b_{ij}=b_{ji}\) 且 \(\mathrm{Q}(\sum_{i}a_{i}e_{i})=\sum_{\{i,j\}}b_{ij}a_{i}a_{j}\) 对 \(a_{i}\in A\) 成立（命题 2）。由于诸元素 \(1\otimes_{A}e_{i}\) 组成

A'-模 A' ⊗\_A E 的基，故在该模上取下式便定义了该模上的一个二次型 Q'

\[
\mathrm{Q} ^ {\prime} (\sum_ {i} a _ {i} ^ {\prime} \otimes e _ {i}) = \sum_ {\{i, j \}} a _ {i} ^ {\prime} a _ {j} ^ {\prime} h (b _ {i j})\tag{20}
\]

其中 \(a_i' \in A'\)；于是，对 \(x = \sum_{i} a_i e_i \in E\)

\[
Q ^ {\prime} (1 \otimes x) = Q ^ {\prime} (\sum_ {i} h (a _ {i}) \otimes e _ {i}) = \sum_ {\{i, j \}} h (a _ {i}) h (a _ {j}) h (b _ {i j}) = h (Q (x)),\tag{21}
\]

这就证明了这种情形下 Q' 的存在性。

现在转到一般情形。设 \((x_{i})_{i\in I}\) 是 E 的一个生成组，\(\mathbf{A}^{(1)}\) 是 I 的元素的形式线性组合模（第 II 章，§ 1，第 8 小节），\((e_{i})_{i\in I}\) 是 \(\mathbf{A}^{(1)}\) 的典范基。线性映射 \(u\) （\(\mathbf{A}^{(1)}\) 到 E 中的、由 \(u(e_{i}) = x_{i}\) 定义的）是满射，因为诸元素 \(x_{i}\) 生成 E。由此（第 III 章，第 2 版，附录 II，第 5 小节，命题 4）可知映射 \(1\otimes u\) （\(\mathbf{A}'\otimes\mathbf{A}^{(1)}\) 到 \(\mathbf{A}'\otimes\mathbf{E}\) 中的）是满射，且它的核 P' 由形如 \(1\otimes p\) 且 \(u(p)=0\) 的元素生成。于是设 \(\mathrm{Q}_{1}^{\prime}\) 是 \(\mathbf{A}'\otimes_{\mathbf{A}}\mathbf{A}^{(1)}\) 上把二次型 \(\mathrm{Q}_{1}=\mathrm{Q}\circ u\) （\(\mathbf{A}^{(1)}\) 上的）加以扩张得到的二次型。若 \(p\) 是 \(\mathbf{A}^{(1)}\) 中满足 \(u(p)=0\) 的元素，则有 \(\mathrm{Q}_{1}^{\prime}(1\otimes p)=h(\mathrm{Q}_{1}(p))=0\)，且（以 \(\Phi_{1}^{\prime}\) 表示与 \(\mathrm{Q}_{1}^{\prime}\) 相伴的双线性型）\(\Phi_{1}^{\prime}(1\otimes p,1\otimes x)=h(\Phi(u(p),u(x)))=0\)（对一切 \(x\in\mathbf{A}^{(1)}\)）。因此，若 \(u(p)=0\)（ \(p\in\mathbf{A}^{(1)}\) ），则 \(1\otimes p\) 属于二次模 \(\mathbf{A}'\otimes_{\mathbf{A}}\mathbf{A}^{(1)}\) 的核，从而如上所述，存在 \(\mathbf{A}'\otimes_{\mathbf{A}}\mathbf{E}\) 上的一个二次型 Q' 使 \(\mathrm{Q}_{1}^{\prime}=\mathrm{Q}^{\prime}\circ(1\otimes u)\)。由于 \(u\) 是满射，可见 Q' 满足条件 (17)。证毕。

其存在性与唯一性由命题 3 保证的二次型 Q′ 称为 Q 到 A′ ⊗\_A E 的（关于 h 的）扩张。也称 Q' 是由 Q 经标量环扩张得到的。

习题. — 1) 设 A 是一个体，E 是 A 上的一个左向量空间，\(\Phi\) 是 E 上的一个半双线性型（关于 A 的某个反自同构 J）。假设 \(\Phi\) 的秩（有限或无限）\(\geqslant 2\)，且关系 \(\Phi(x, y) = 0\) 与 \(\Phi(y, x) = 0\) 等价。

\begin{enumerate}
\def\labelenumi{\alph{enumi})}
\item
  证明：存在 A 中的 \(\lambda\neq0\) 使 \(\Phi(y,x)=\lambda(\Phi(x,y))^{\mathfrak{J}}\)（利用 § 1 习题 8）。
\item
  证明：存在 \(\alpha\in A\)，使半双线性型 \(\Phi\alpha\)（关于反自同构 \(\xi\to\alpha^{-1}\xi^{J}\alpha\)）是埃尔米特型或交错型。
\end{enumerate}

（首先注意有 \(\xi^{J^2} = \lambda^{-1}\xi\lambda^{-J}\) 且 \(\lambda\lambda^J = \lambda^J\lambda = 1\)。然后区分两种情形：是 \(\xi + \xi^J\lambda^{-1} = 0\) 对一切 \(\xi \in A\) 成立，还是不然；在第二种情形，证明每个 \(\neq 0\) 的形如 \(\alpha = \xi + \xi^J\lambda^{-1}\) 的元素都满足要求。）

\begin{enumerate}
\def\labelenumi{\arabic{enumi})}
\setcounter{enumi}{1}
\tightlist
\item
  设 \(\Phi\) 是体 A 上有限维向量空间 E 上的 \(\varepsilon\)-埃尔米特半双线性型。
\end{enumerate}

\begin{enumerate}
\def\labelenumi{\alph{enumi})}
\item
  证明：对 E 的每个向量子空间 M，有 \((\mathbf{M}^{\mathbf{0}})^{\mathbf{0}} = \mathbf{M} + \mathbf{E}^{\mathbf{0}}\) 且 \(\dim M^{\mathbf{0}} + \dim M = \dim E + \dim (M \cap E^{\mathbf{0}})\)。
\item
  若 \(M_{1}\)、\(M_{2}\) 是 E 的两个向量子空间，证明有 dim \((\mathbf{M}_{1} \cap \mathbf{M}_{2}^{0}) + \dim (\mathbf{M}_{2} + \mathbf{M}_{1}^{0}) = \dim \mathbf{E} + \dim (\mathbf{M}_{1} \cap \mathbf{E}^{0})\)（考察 E 到 \(E/E^{0}\) 的典范映射）。
\end{enumerate}

\begin{enumerate}
\def\labelenumi{\arabic{enumi})}
\setcounter{enumi}{2}
\tightlist
\item
  设 K 是一个交换环，\(f\) 是 K{[}X{]} 中的一个首一多项式，次数为 \(n \geqslant 1\)；设 A 是商代数 K{[}X{]}/(f)，它以 1, \(\xi, \xi^2, \ldots, \xi^{n-1}\) 为 K 上的基（第 IV 章，§ 1，第 5 小节，命题 4）。证明：一个 A-模 E 的数据等价于一个 K-模 \(E_0\) 与 K-自同态 \(j\) （\(E_0\) 的、满足 \(f(j) = 0\) 的）的数据。对每个
\end{enumerate}

\(u \in \mathrm{E}^*\)，令 \(u(x) = \sum_{k=0} u_k(x)\xi^k\)；证明：若 \(\alpha_0 = f(0)\) 在 K 中可逆，则映射 \(u \to u_0\) 是 \(\mathbf{E}^*\) 到 \((\mathbf{E}_0)^*\) 上的一个 K-同构，并写出其逆同构。

¶ 4) a) 设 A 是一个环（交换或不交换），σ 是 A 的一个自同构，使得存在可逆元素 γ ∈ A 满足 γ\^{}σ = γ，且对一切 ξ ∈ A 有 ξ\^{}σ² = γξγ\^{}-1。设 B 是具有两个元素的基 (e₁, e₂) 的左 A-模；证明：在 B 中取下式作为乘法，便在 B 上定义了一个环结构

\[
(\xi e _ {1} + \eta e _ {2}) (\xi^ {\prime} e _ {1} + \eta^ {\prime} e _ {2}) = (\xi \xi^ {\prime} + \eta \eta^ {\prime \sigma} \gamma) e _ {1} + (\eta \xi^ {\prime \sigma} + \xi \eta^ {\prime}) e _ {2}.
\]

对这个环结构，\(e_1\) 是单位元（与 A 的单位元 1 视为同一）；若令 \(e_2 = \rho\)，则 \(\rho^2 = \gamma\) 且 \(\rho\xi = \xi^\sigma\rho\)（对一切 \(\xi \in A\)）。此外，B 是一个右 A-模，1 与 \(\rho\) 组成它的基。若 A 是一个体，则 B 是体的必要充分条件是 \(\gamma\) 不是形如 \(\lambda^\sigma\lambda\)（其中 \(\lambda \in A\)）的元素。（参见第 VIII 章，§ 12，习题 8。）

\begin{enumerate}
\def\labelenumi{\alph{enumi})}
\setcounter{enumi}{1}
\tightlist
\item
  设 J 是 A 的一个对合反自同构。假设存在可逆元素 \(\delta\in\mathsf{A}\) 满足下列条件：
\end{enumerate}

\[
\delta^ {J} = \delta , \quad \delta^ {\sigma} \delta = \gamma \gamma^ {J}, \quad (\xi^ {J}) ^ {\sigma} = \delta (\xi^ {\sigma}) ^ {J} \delta^ {- 1} \quad \text { for   all } \quad \xi \in A.\tag{1}
\]

证明：这时可把 J 扩张为 B 的一个对合反自同构（仍记作 J），办法是令

\[
(\xi + \eta \rho) ^ {J} = \xi^ {J} + \gamma^ {- 1} \delta^ {\sigma} (\eta^ {J}) ^ {\sigma} \rho .\tag{2}
\]

此外，若 A 与 B 都是体，且 \(\sigma\) 不是内自同构，证明条件 (1) 是使 J 可扩张为 B 的对合反自同构的必要条件（令 \(\rho^{J} = \alpha + \beta\rho\)，其中 \(\alpha, \beta\) 属于 A，证明必有 \(\alpha = 0\)，其法是写出条件 \((\rho\xi)^{J} = \xi^{J}\rho^{J}\)（对一切 \(\xi \in A\)））。

\begin{enumerate}
\def\labelenumi{\alph{enumi})}
\setcounter{enumi}{2}
\tightlist
\item
  假设条件 (1) 成立，且对合反自同构 J 已由 (2) 扩张到 B 上。设 F 是一个单式 B-模，E 是把 F 的标量环限制到 A 而得到的单式 A-模；映射 \(j: x \to \rho x\) 于是是 E 到其自身上的一个半线性双射，它关于 A 的自同构 \(\sigma\)，且满足 \(j^{2}(x) = \gamma x\)。设 \(\Phi\) 是 F 上的一个埃尔米特型（关于 J）；对 \(x \in E\)、\(y \in E\)，令 \(\Phi(x, y) = \Phi_{1}(x, y) + \Phi_{2}(x, y)\rho\)，其中 \(\Phi_{1}(x, y) \in A\)、\(\Phi_{2}(x, y) \in A\)。证明 \(\Phi_{1}\) 是 E 上的一个埃尔米特型（关于 J），且满足
\end{enumerate}

\[
\Phi_ {1} (j (x), j (y)) = (\Phi_ {1} (x, y)) ^ {\sigma} \delta ;
\]

我们有 \(\Phi_{2}(x,y)=\Phi_{1}(x,j(y))\delta^{-1}\)，\(\Phi_{2}\) 是 E 上关于反自同构（一般非对合）\(\xi\to(\xi^{\mathrm{J}})^{\sigma}\) 的半双线性型，且满足

\[
\Phi_ {2} (j (x), j (y)) = (\Phi_ {2} (x, y)) ^ {\sigma} \delta^ {\sigma}
\]

and

\[
\Phi_ {2} (y, x) = \gamma^ {\mathrm{J}} \delta^ {- 1} ((\Phi_ {2} (x, y)) ^ {\mathrm{J}}) ^ {\sigma}.
\]

逆命题。要使 \(\Phi\) 非退化，必须且只需 \(\Phi_{1}\) 或 \(\Phi_{2}\) 非退化。B 是交换环 K 上对应于 K 的元素对 \((\alpha, \beta)\)（其中 \(\beta\) 可逆）的四元数代数、A 是 B 的子代数 K + Ku、J 与 \(\sigma\) 是映射 \(\xi \to \bar{\xi}\)（第 II 章，§ 7，第 8 小节）的特殊情形。

\begin{enumerate}
\def\labelenumi{\alph{enumi})}
\setcounter{enumi}{3}
\tightlist
\item
  设 u 是 A-模 E 的一个自同构。要使 u 是 B-模 F 的保持型 \(\Phi\) 不变的自同构，必须且只需 u 满足下列三个条件中的任意两个：
\end{enumerate}

\(1^{\circ} u\) 保持 \(\Phi_{1}\) 不变；

\(2^{o} u\) 保持 \(\Phi_{2}\) 不变；

\(3^{o} u\) 与 \(j\) 交换。

\begin{enumerate}
\def\labelenumi{\arabic{enumi})}
\setcounter{enumi}{4}
\tightlist
\item
  设 A 是一个交换环，E 是一个 A-模，\((x_{i})_{i\in I}\) 是 E 的一个生成组；设 R 是模 \(\mathrm{A}^{(\mathrm{I})}\) 中由元素 \((y_{i})_{i\in\mathrm{I}}\)（满足 \(\sum_{i}y_{i}x_{i}=0\)）组成的子模，\((a_{\lambda})_{\lambda\in\mathrm{L}}\) 是 R 的一个生成组（其中 \(a_{\lambda}=(a_{\lambda i})_{i\in\mathrm{I}}\)）。设 \((b_{ij})\) 是 A 的元素的一个族 \((i\in\mathrm{I},j\in\mathrm{I})\)。要存在一个二次型 Q 使 \(\mathrm{Q}(x_{i})=b_{ii}\) 且 \(\Phi(x_{i},x_{j})=b_{ij}\)（对 \(i\neq j\)）（ \(\Phi\) 表示与 Q 相伴的双线性型），必须且只需对 I 中一切 i, j 有 \(b_{ij}=b_{ji}\)，且对一切 \(\lambda\in L\) 与 \(i\in I\) 有
\end{enumerate}

\[
\sum_ {\{i, j \}} b _ {i j} a _ {\lambda i} a _ {\lambda j} = \sum_ {j \neq i} b _ {i j} a _ {\lambda j} + 2 b _ {i i} a _ {\lambda i} = 0;\tag{3}
\]

于是有 \(\mathrm{Q}(\sum_{i}a_{i}x_{i})=\sum_{\{i,j\}}b_{ij}a_{i}a_{j}\)。由此推出第 4 小节命题 3 的一个新证明。（注意 \(x_{i}^{\prime}=1\otimes x_{i}\) 组成 \(A^{\prime}\otimes_{A}E\) 的一个生成组，且 \(A^{\prime}\)-模 \(A^{\prime}\otimes_{A}E\) 同构于 \(A^{\prime(I)}/R^{\prime}\)，其中 \(A^{\prime(I)}\) 与 \(A^{\prime}\otimes_{A}A^{(I)}\) 视为同一，而 \(R^{\prime}\) 由 R 在 \(A^{(I)}\) 到 \(A^{\prime(I)}\) 中的典范映射下的像生成。）

\begin{enumerate}
\def\labelenumi{\arabic{enumi})}
\setcounter{enumi}{5}
\tightlist
\item
  设 A 是特征 2 的交换环，E 是自由 A-模，\(\mathcal{A}\)（相应地 \(\mathcal{S},\mathcal{Q}\)）是 E 上交错（相应地对称双线性、二次）型所成的 A-模。我们有 \(\mathcal{A}\subset\mathcal{S}\)；再定义线性映射 \(\omega\) （\(\mathcal{S}\) 到 \(\mathcal{Q}\) 中的）与线性映射 \(\theta\) （\(\mathcal{Q}\) 到 \(\mathcal{A}\) 中的）如下：对每个双线性型 \(\Phi\in\mathcal{S}\)，\(\omega(\Phi)\) 是二次型 \(x\to\Phi(x,x)\)；对每个二次型 \(Q\in\mathcal{Q}\)，\(\theta(Q)\) 是与 Q 相伴的双线性型，它是交错的。证明 \(\omega(0)=\mathcal{A}\)、\(\theta(\mathcal{Q})=\mathcal{A}\) 且 \(\bar{\theta}(0)=\omega(\mathcal{S})\)。
\end{enumerate}

¶ 7) 设 A 是一个交换环，E、F 是两个 A-模。称 E 到 F 中的映射 Q 为二次映射，如果它满足下列条件：\(1^{\circ} \mathrm{Q}(\alpha x) = \alpha^{2} \mathrm{Q}(x)\)（对 \(\alpha \in \mathrm{A}\)、\(x \in \mathrm{E}\)）；\(2^{\circ}\) 映射 \((x, y) \to \mathrm{Q}(x + y) - \mathrm{Q}(x) - \mathrm{Q}(y)\) （\(\mathrm{E} \times \mathrm{E}\) 到 F 中的）是双线性的。若 \(f\) 是 A-模 \(\mathrm{E}_{1}\) 到 E 中的线性映射，则 \(\mathrm{Q} \circ f\) 是 \(\mathrm{E}_{1}\) 到 F 中的二次映射。

\begin{enumerate}
\def\labelenumi{\alph{enumi})}
\tightlist
\item
  设 E 是一个 A-模，\(\mathrm{A}^{(\mathrm{E})}\) 是 E 的元素以 A 为系数的形式线性组合模（第 II 章，§ 1，第 8 小节），且对每个 \(x \in \mathrm{E}\)，设 \(\varepsilon_x\) 是 \(\mathrm{A}^{(\mathrm{E})}\) 的典范基中相应的元素。设 \(\Gamma^2(\mathrm{E})\) 是 \(\mathrm{A}^{(\mathrm{E})} \times (\mathrm{E} \otimes_{\mathbf{A}} \mathrm{E})\) 关于由形如 ( \(\varepsilon_{x+y} - \varepsilon_x - \varepsilon_y, -x \otimes y\) ) 与 ( \(\varepsilon_{\lambda x} - \lambda^2\varepsilon_x, 0\) )（对 \(x \in \mathrm{E}\)、\(y \in \mathrm{E}\)、\(\lambda \in \mathrm{A}\)）的元素生成的子模 R 的商。对每个 \(x \in \mathrm{E}\)，令 \(\gamma(x) = \varphi(\varepsilon_x, 0)\)，其中 \(\varphi\) 表示 \(\mathrm{A}^{(\mathrm{E})} \times (\mathrm{E} \otimes \mathrm{E})\) 到 \(\Gamma^2(\mathrm{E})\) 上的典范映射；称 \(\gamma\) 为 E 到 \(\Gamma^2(\mathrm{E})\) 中的典范映射。证明 \(\gamma\) 是 E 到 \(\Gamma^2(\mathrm{E})\) 的二次映射，且对每个 E 到 A-模 F 中的二次映射 Q，存在唯一的线性映射 \(q\) （\(\Gamma^2(\mathrm{E})\) 到 F 中的）使 \(Q = q \circ \gamma\)（换言之，\((\Gamma^2(\mathrm{E}), \gamma)\) 是一个泛映射问题的解；参见集合论，第 IV 章，§ 3）。
\end{enumerate}

对每对 A-模 E, E′ 以及 E 到 E′ 中的每个线性映射 f，证明：若 γ′ 表示 E′ 到 Γ²(E′) 中的典范映射，则存在唯一的 Γ²(E) 到 Γ²(E') 中的线性映射 \(\bar{f}\) 使 γ'∘f = \(\bar{f} \circ \gamma\)。

\begin{enumerate}
\def\labelenumi{\alph{enumi})}
\setcounter{enumi}{1}
\item
  假设 E 是两个子模 M, N 的直和。定义 \(\Gamma^{2}(\mathrm{E})\) 到模 \(\Gamma^{2}(\mathrm{M})\)、\(\Gamma^{2}(\mathrm{N})\) 与 \(\mathbf{M} \otimes \mathbf{N}\) 的直和上的一个典范同构（证明此直和是与 \(\Gamma^{2}(\mathrm{E})\) 相同的泛映射问题的解）。
\item
  设 F 是 E 的一个子模，j 是 F 到 E 中的典范单射。定义 \(\Gamma^{2}(\mathrm{E}/\mathrm{F})\) 到
\end{enumerate}

\[
\Gamma^ {2} (\mathrm{E}) / (\bar {j} (\Gamma^ {2} (\mathrm{F})) + \psi (\mathrm{E} \times \mathrm{F})),
\]

其中 \(\psi(x, y) = \varphi(0, x \otimes j(y))\)（对 \(x \in \mathrm{E}, y \in \mathrm{F}\)）。（同一方法。）

\begin{enumerate}
\def\labelenumi{\alph{enumi})}
\setcounter{enumi}{3}
\item
  设 A' 是一个交换环，h 是 A 到 A' 中的一个同态。定义 \(\Gamma^{2}(\mathbf{A}'\otimes_{\mathbf{A}}\mathbf{E})\) 到 \(\mathbf{A}'\otimes_{\mathbf{A}}\symbf{\Gamma}^{2}(\mathbf{E})\) 上的一个典范同构（同一方法）。
\item
  存在唯一的线性映射 \(s\) （\(\Gamma^{2}(\mathrm{E})\) 到 \(\mathrm{E}\otimes\mathrm{E}\) 中的），使 \(s(\gamma(x))=x\otimes x\) 对一切 \(x\in\mathrm{E}\) 成立；证明若 \(\mathrm{E}\) 是自由模，则 s 是到 E 上 2 阶对称张量所成的子模上的一个同构。
\item
  假设 A = Z 且 E 是 n 阶有限循环群。证明：若 n 为奇数，\(\Gamma^{2}(E)\) 是 n 阶循环群；若 n 为偶数，则是 2n 阶循环群。（首先注意，若 a 是 E 的一个生成元，则 \(\gamma(a)\) 是 \(\Gamma^{2}(E)\) 的生成元，且对每个整数 h 有 \(\gamma(-ha) = \gamma(ha)\)；由此推出，若 n 为奇数，则 \(n\gamma(a) = 0\)（取 \(h = (n - 1)/2\)）；同样证明若 n 为偶数则 \(2n\gamma(a) = 0\)。最后证明：若 n 为奇数（相应地偶数），存在 E 到 n（相应地 2n）阶循环群中的把 a 映到该群生成元上的二次映射 Q。）
\end{enumerate}

\begin{enumerate}
\def\labelenumi{\arabic{enumi})}
\setcounter{enumi}{7}
\tightlist
\item
  设 A 是一个域，E、F 是 A 上的两个向量空间。设 g 是 E 到 F 中的一个映射，使得存在 E × E 到 F 中的三个映射 a, b, c，满足恒等式
\end{enumerate}

\[
g (\lambda x + \mu y) = \lambda^ {2} a (x, y) + \lambda \mu b (x, y) + \mu^ {2} c (x, y)\tag{4}
\]

对 E 中一切 x, y 与 A 中 λ, μ 成立。

\begin{enumerate}
\def\labelenumi{\alph{enumi})}
\tightlist
\item
  证明有 \(a(x, y) = g(x)\)、\(c(x, y) = g(y)\)、\(b(y, x) = b(x, y)\) 及 \(b(\lambda x, y) = \lambda b(x, y)\)；此外，若令
\end{enumerate}

\[
d (x, y, z) = b (x + y, z) - b (x, z) - b (y, z)
\]

证明有 \(d(x,y,z)=d(y,z,x)=d(z,x,y)\)，并由此推出 \(d(\lambda x,\mu y,\nu z)=\lambda\mu\nu d(x,y,z)\)。

\begin{enumerate}
\def\labelenumi{\alph{enumi})}
\setcounter{enumi}{1}
\tightlist
\item
  由 a) 推出：若 \(\mathbf{A} \neq \mathbf{F}_2\)，则必有 \(d(x, y, z) = 0\)，从而 \(g\) 是二次映射（习题 7）。反之，若 \(\mathbf{A} = \mathbf{F}_2\) 且 dim \(\mathrm{E} \geqslant 3\)，证明存在 E 到 A 中的映射 \(g\)，满足 (4) 且使 \(d(x, y, z)\) 不恒为零。
\end{enumerate}

\begin{enumerate}
\def\labelenumi{\arabic{enumi})}
\setcounter{enumi}{8}
\tightlist
\item
  设 A 是一个交换环，E 是具有 n 个元素的基的 A-模，Φ 是 E 上的对称双线性型。设 e 是 \(\bigwedge^{n}\) E 中组成该模的基的一个元素，\(\Delta\) 是 \(\Phi\) 关于 \(e\) 的判别式，\(\varphi_{p}\) 是 \(\bigwedge^{p}\) E 到 \(\bigwedge^{n-p}\) E* 上的典范同构（关于 \(e\)）（第 III 章，§8，第 5 小节）。对 \(x \in \bigwedge^{p}\) E，令 \(d_{(p)}(x) = \bar{\varphi}_{n-p}^{-1}(d_{\Phi(p)}(x))\)；证明映射 \(d_{(p)}\) （\(\bigwedge^{p}\) E 到 \(\bigwedge^{n-p}\) E 中的）具有下列性质：
\end{enumerate}

\begin{enumerate}
\def\labelenumi{\alph{enumi})}
\item
  对每个 \(x \in \wedge \mathrm{E}\)，有 \(d_{(n - p)}(d_{(p)}(x)) = (-1)^{p(n - p)}\Delta x\)。
\item
  对每对元素 \(x, y\) （属于 \(\wedge\) E），有
\end{enumerate}

\[
x \wedge d _ {(p)} (y) = \Phi_ {(p)} (x, y) e \quad \text { and } \quad \Phi_ {(n - p)} (d _ {(p)} (x), d _ {(p)} (y)) = \Delta \Phi_ {(p)} (x, y).
\]

\begin{enumerate}
\def\labelenumi{\alph{enumi})}
\setcounter{enumi}{2}
\item
  进一步假设 A 是一个体且 \(\Phi\) 非退化。那么，若 \(x\) 是对应于 E 的子空间 F 的一个可分解 \(p\)-向量 \(\neq 0\)（第 III 章，§ 7，第 3 小节），则 \(d_{(p)}(x)\) 是一个可分解 \((n-p)\)-向量 \(\neq 0\)，对应于与 F 正交的子空间 F \(^{0}\)。
\item
  将前面的结果推广到如下情形：（A 为交换环，）\(\Phi\) 是 \(\varepsilon\)-埃尔米特半双线性型（关于 A 的一个对合自同构 \(J \neq 1\)）。
\end{enumerate}

\begin{enumerate}
\def\labelenumi{\arabic{enumi})}
\setcounter{enumi}{9}
\tightlist
\item
  \begin{enumerate}
  \def\labelenumii{\alph{enumii})}
  \tightlist
  \item
    设 A 是一个交换环，E 是具有 3 个元素的基的 A-模，\(\Phi\) 是 E 上的对称双线性型。沿用习题 9 的记号，对 E 的任意两个元素 \(x\)、\(y\)，令 \(x \overline{\wedge} y = d_{(2)}(x \wedge y)\)，并称此元素为向量积
  \end{enumerate}
\end{enumerate}

\[
\stackrel {{3}} {\wedge} \mathrm{E})
\]

即 \(x\) 与 \(y\) 的向量积（关于 \(\Phi\) 以及 \(e\) （\(\wedge E\) 的基））。证明 \((x, y) \to x \overline{\wedge} y\) 是 \(E \times E\) 到 \(E\) 中的交错双线性映射，且 \(x \overline{\wedge} y\) 与 \(x\) 以及 \(y\) 都正交。

\begin{enumerate}
\def\labelenumi{\alph{enumi})}
\setcounter{enumi}{1}
\tightlist
\item
  设 \(\alpha, \beta\) 是 A 的两个可逆元素，B 是与偶 \((\alpha, \beta)\) 对应的 A 上的四元数代数（第 II 章，§ 7，第 8 小节），1, u, v, w 是 B 在 A 上的典范基；设 E 是以 u, v, w 为基的 B 的子模。证明：若 x, y 是属于 E 的两个四元数，则有
\end{enumerate}

\[
x y = \Phi (x, y) + x \overline {{{{\wedge}}}} y
\]

其中 \(\Phi\) 是 E 上的对称双线性型，与 \(\Phi\) 相伴的线性映射是双射，而 \(x\overline{\wedge}y\) 是 x 与 y 关于型 \(\Phi\) 以及基 \(\alpha^{-1}\beta^{-1}u\wedge v\wedge w\) （\(\bigwedge^{3}\mathbf{E}\) 的）的向量积。

\begin{enumerate}
\def\labelenumi{\arabic{enumi})}
\setcounter{enumi}{10}
\tightlist
\item
  设 \(\Phi\) 是有限维向量空间 E 上的非退化 \(\varepsilon\)-埃尔米特半双线性型。若两个向量子空间 M、N° 中的一个包含另一个，则称 E 的向量子空间 M （关于 \(\Phi\)）与向量子空间 N 弱正交。
\end{enumerate}

\begin{enumerate}
\def\labelenumi{\alph{enumi})}
\item
  证明关系“M 与 N 弱正交”是对称的。
\item
  若 M 与 N 弱正交，则 \(M^{0}\) 与 \(N^{0}\) 弱正交。
\item
  若 M 与 N 弱正交且 M ∩ N = \{0\}，则 M 与 N 正交。
\end{enumerate}

